\section{Results}

\subsection{Physical cohort}

All 10 site-specific synthetic head-and-neck cases completed the physical transport workflow. Peripheral target coverage was stable across the cohort, with mean PTV $D_{95}$ fixed at 3.50~Gy by design, while lattice selectivity varied more broadly. The median physical PVDR was approximately 1.80, and the physical cohort retained measurable spread in peri-tumour spill and OAR-adjacent dose. These distributions established the physical baseline against which biological reinterpretation was tested (Figure~\ref{fig:physical_cohort}).

\begin{figure}[H]
\centering
\includegraphics[width=0.92\textwidth]{figureR1_physical_cohort_summary.png}
\caption{\textbf{Physical cohort summary.} Across 10 synthetic head-and-neck cases, the TOPAS-derived photon lattice plans achieved stable peripheral target coverage but retained broad variation in peak-valley contrast, spill, and OAR-adjacent physical dose.}
\label{fig:physical_cohort}
\end{figure}

\subsection{Biological reinterpretation of the 10-case cohort}

Applying the biology-informed post-processing substantially changed cohort interpretation. Relative to the physical-only view, cohort-mean PTV $D_{95}$ increased from 3.50~Gy to 17.14~Gy in the no-sink analysis and 15.73~Gy with the anatomical sink included. In contrast, mean PVDR fell from 1.817 physically to 1.478 without sink uptake and 1.436 with sink uptake, indicating biological filling-in of the physical peak-valley pattern. Peri-GTV 0--5~mm mean increased from 7.219~Gy physically to 20.616~Gy without sink uptake and 19.929~Gy with sink uptake. Brainstem $D_{2}$ increased from 11.454~Gy physically to 28.399~Gy without sink uptake and 27.951~Gy with sink uptake. Right parotid mean increased from 5.875~Gy physically to 18.821~Gy without sink uptake and 17.939~Gy with sink uptake (Table~\ref{tab:cohort_mode_summary}).

\begin{table}[H]
\centering
\caption{\textbf{Cohort-mean endpoint summary across comparison modes.}}
\label{tab:cohort_mode_summary}
\small
\begin{tabular}{p{3.9cm}ccc}
\toprule
Endpoint & Physical-only & No-sink & Anatomical sink \\
\midrule
PTV $D_{95}$ (Gy) & 3.500 & 17.144 & 15.731 \\
PVDR & 1.817 & 1.478 & 1.436 \\
Peri-GTV 0--5 mm mean (Gy) & 7.219 & 20.616 & 19.929 \\
Peri-GTV 5--15 mm mean (Gy) & 5.923 & 16.900 & 16.492 \\
Brainstem $D_{2}$ (Gy) & 11.454 & 28.399 & 27.951 \\
Right parotid mean (Gy) & 5.875 & 18.821 & 17.939 \\
\bottomrule
\end{tabular}
\end{table}


Biological reinterpretation changed the final ordering of 8 of 10 cases relative to physical-only ranking. The best-ranked physical case was the oral tongue/floor-of-mouth horseshoe case, but it dropped from rank 1 physically to rank 3 with the anatomical sink included. In contrast, the right sinonasal/maxillary bulky ellipsoid improved from physical rank 2 to biological rank 1. Five cases acquired biology-only OAR reinterpretations, all involving brainstem $D_{2}$. The full 10-case rank-change table is shown in Table~\ref{tab:rank_changes}.

\begin{figure}[H]
\centering
\includegraphics[width=0.92\textwidth]{figureR2_biological_reinterpretation.png}
\caption{\textbf{Biological reinterpretation of the site-specific cohort.} The same physical cohort was rescored under biology-informed post-processing with and without anatomical vascular uptake. Plan ordering changed in most cases, and spill and OAR burden rose substantially relative to the physical-only view.}
\label{fig:bio_reinterpretation}
\end{figure}

\begin{table}[H]
\centering
\caption{\textbf{Rank changes across the 10-case cohort after biology-informed reinterpretation.} Negative shift values indicate improvement relative to the physical-only rank.}
\label{tab:rank_changes}
\scriptsize
\begin{tabular}{p{0.7cm}p{3.8cm}p{0.8cm}p{0.8cm}p{0.9cm}p{1.5cm}p{2.4cm}}
\toprule
Case & Site / tumour model & Physical rank & No-sink rank & Anatomical sink rank & Sink vs physical shift & Biology-only reinterpretation \\
\midrule
07 & Oral tongue / floor-of-mouth horseshoe & 1 & 2 & 3 & +2 & None \\
01 & Right sinonasal / maxillary bulky ellipsoid & 2 & 1 & 1 & -1 & None \\
02 & Left maxillary crescent abutting orbit & 3 & 3 & 2 & -1 & None \\
04 & Laryngo-hypopharyngeal elongated cylinder & 4 & 4 & 4 & 0 & Brainstem $D_{2}$ \\
05 & Parapharyngeal / prestyloid deep-space mass & 5 & 5 & 6 & +1 & Brainstem $D_{2}$ \\
09 & Deep parotid / infratemporal bulky mass & 6 & 9 & 9 & +3 & Brainstem $D_{2}$ \\
06 & Buccal mucosa / cheek infiltrative crescent & 7 & 6 & 5 & -2 & None \\
08 & Bulky nodal conglomerate with central necrosis & 8 & 7 & 7 & -1 & None \\
03 & Base-of-tongue / oropharynx bilobed mass & 9 & 8 & 8 & -1 & Brainstem $D_{2}$ \\
10 & Composite oropharynx + upper-neck bulky tumour & 10 & 10 & 10 & 0 & Brainstem $D_{2}$ \\
\bottomrule
\end{tabular}
\end{table}


\subsection{Uncertainty analysis}

Repeat-run uncertainty was evaluated on a six-case subset using three seeds and three history scales, giving nine repeated transport combinations. Endpoint-level sink effects remained visible after uncertainty was introduced. Across the evaluated case-endpoint set, 40 of 42 sink-driven endpoint deltas exceeded the combined Monte Carlo plus uptake-sensitivity 95\% noise band. The strongest repeated signals were observed for PVDR, peri-GTV spill, and right parotid mean. However, no case showed a noise-qualified sink-driven rank change in this subset, indicating that the sink effect was more stable as an endpoint-level modifier than as a rank-reversal mechanism (Figure~\ref{fig:uncertainty_robustness}).

\begin{figure}[H]
\centering
\includegraphics[width=0.92\textwidth]{figureR3_uncertainty_robustness.png}
\caption{\textbf{Uncertainty robustness.} Repeated transport and biology solves preserved most sink-driven endpoint shifts beyond the combined 95\% noise band, but did not produce stable sink-driven rank reversals in the repeated six-case subset.}
\label{fig:uncertainty_robustness}
\end{figure}

\subsection{Vascular-sink falsification}

The anatomical sink remained distinguishable from stronger surrogate controls. In the repeated subset, 222 of 288 true-versus-surrogate endpoint differences exceeded the combined 95\% noise band. The strongest repeated separations were observed for the no-sink comparator (42/48 endpoints), uniform body washout (42/48), and local vessel dropout near the tumour (41/48). Superior-inferior and anterior-posterior flips remained weaker but still measurable, with 32/48 and 31/48 endpoint differences above the same threshold, respectively. Only 2 of 36 true-versus-surrogate rank comparisons were noise-qualified, and both occurred against the mass-matched uniform washout control. Thus, the anatomical sink was reproducibly different from generic and spatially perturbed alternatives at the endpoint level, but it remained a weak rank-flipping signal (Table~\ref{tab:falsification_summary}; Figure~\ref{fig:falsification}).

\begin{table}[H]
\centering
\caption{\textbf{Summary of stronger falsification results in the repeated subset.}}
\label{tab:falsification_summary}
\small
\begin{tabular}{p{5.6cm}cc}
\toprule
Comparator & Endpoints above 95\% noise band & Noise-qualified rank changes \\
\midrule
No sink & 42/48 & 0/6 \\
Uniform body washout (mass matched) & 42/48 & 2/6 \\
Local vessel dropout within 20 mm of tumour & 41/48 & 0/6 \\
Randomized anatomy-constrained displacement & 34/48 & 0/6 \\
Superior-inferior flip & 32/48 & 0/6 \\
Anterior-posterior flip & 31/48 & 0/6 \\
\bottomrule
\end{tabular}
\end{table}


\begin{figure}[H]
\centering
\includegraphics[width=0.92\textwidth]{figureR4_stronger_falsification_controls.png}
\caption{\textbf{Stronger spatial falsification controls.} The anatomical sink remained distinguishable from multiple perturbed or redistributed washout fields, especially from local vessel dropout and uniform body washout, even after repeat-run uncertainty was included.}
\label{fig:falsification}
\end{figure}

\subsection{Summary}

Taken together, the cohort results support three points. First, the physical TOPAS-derived photon lattice plans were reproducibly generated across 10 anatomically distinct synthetic H\&N cases. Second, biology-informed post-processing materially changed case interpretation, including rank order and OAR burden. Third, the vascular-sink term behaved as a measurable endpoint-level modifier that remained distinguishable from several stronger surrogate controls, although it did not produce broad noise-qualified rank reversals.
